% GENERATED FILE - do not edit.
% Built from resume-sde.tex by tools/flatten.py; run `make standalone`.
% Self-contained: upload this one file to Overleaf and compile.
%-------------------------
% Resume in LaTeX
% Variant: SOFTWARE ENGINEERING / BACKEND internships
% ATS-optimised. Compile with pdfLaTeX.
%
% Shared styling lives in preamble.tex; shared sections live in sections/.
% Only SKILLS and PROJECTS differ between this file and resume-aiml.tex.
%------------------------

\documentclass[letterpaper,10pt]{article}

% >>> inlined from preamble.tex (generated by tools/flatten.py; edit the original)
%-------------------------
% Shared preamble: styling, fonts and resume macros.
% Both resume-sde.tex and resume-aiml.tex \input this file, so a change to
% spacing or colours here applies to every variant at once.
% Compile with pdfLaTeX.
%------------------------

\input{glyphtounicode}
\pdfgentounicode=1

\usepackage{latexsym}
\usepackage[empty]{fullpage}
\usepackage{titlesec}
\usepackage{marvosym}
\usepackage[usenames,dvipsnames]{color}
\usepackage{verbatim}
\usepackage{enumitem}
\usepackage[hidelinks]{hyperref}
\usepackage{fancyhdr}
\usepackage[english]{babel}
\usepackage{tabularx}
\usepackage{fontawesome5}
\usepackage[normalem]{ulem}
\usepackage{tgheros}
\usepackage[T1]{fontenc}
\usepackage[scale=0.90,lf]{FiraMono}

\definecolor{light-grey}{gray}{0.83}
\definecolor{dark-grey}{gray}{0.3}
\definecolor{text-grey}{gray}{.08}

\renewcommand{\ULdepth}{1.8pt}

\renewcommand*\familydefault{\sfdefault}

\pagestyle{fancy}
\fancyhf{}
\fancyfoot{}
\renewcommand{\headrulewidth}{0pt}
\renewcommand{\footrulewidth}{0pt}

\addtolength{\oddsidemargin}{-0.5in}
\addtolength{\evensidemargin}{0in}
\addtolength{\textwidth}{1in}
\addtolength{\topmargin}{-.55in}
\addtolength{\textheight}{1.1in}

\urlstyle{same}
\raggedbottom
\raggedright
\setlength{\tabcolsep}{0in}

\titleformat{\section}{
    \bfseries \vspace{-2pt} \raggedright \large
}{}{0em}{}[\color{light-grey}{\titlerule[2pt]}\vspace{-4pt}]

\newcommand{\resumeItem}[1]{\item\small{{#1 \vspace{-2pt}}}}

\newcommand{\resumeSubheading}[4]{
  \vspace{-2pt}\item
    \begin{tabular*}{\textwidth}[t]{l@{\extracolsep{\fill}}r}
      \textbf{#1} & {\color{dark-grey}\small #2}\vspace{1pt}\\
      \textit{#3} & {\color{dark-grey}\small #4}\\
    \end{tabular*}\vspace{-5pt}
}

\newcommand{\resumeProjectHeading}[2]{
  \item
    \begin{tabular*}{\textwidth}{l@{\extracolsep{\fill}}r}
      #1 & {\color{dark-grey}\small #2}\\
    \end{tabular*}\vspace{-5pt}
}

\newcommand{\resumeSubItem}[1]{\resumeItem{#1}\vspace{-4pt}}
\renewcommand\labelitemii{$\vcenter{\hbox{\tiny$\bullet$}}$}

\newcommand{\resumeSubHeadingListStart}{\begin{itemize}[leftmargin=0in, label={}]}
\newcommand{\resumeSubHeadingListEnd}{\end{itemize}}
\newcommand{\resumeItemListStart}{\begin{itemize}[topsep=0pt]}
\newcommand{\resumeItemListEnd}{\end{itemize}\vspace{-2pt}}

\color{text-grey}
% <<< end preamble.tex

\begin{document}

% >>> inlined from sections/header.tex (generated by tools/flatten.py; edit the original)
%----------HEADING----------
% Each contact item is wrapped in \mbox so its icon can never be split from
% its text at a line break. The visible link text is the URL itself, because
% applicant tracking systems read the text layer and ignore link annotations.
\begin{center}
    \textbf{\Huge Uttkarsh Thakur} \\ \vspace{4pt}
    \small
    \mbox{\faPhone* \hspace{2pt} +91-6202594338} \hspace{1pt} $|$
    \hspace{1pt} \mbox{\faEnvelope \hspace{2pt} \href{mailto:thakuruttkarsh2603@gmail.com}{thakuruttkarsh2603@gmail.com}} \hspace{1pt} $|$
    \hspace{1pt} \mbox{\faLinkedin \hspace{2pt} \href{https://www.linkedin.com/in/uttkarsh-thakur-76979228a/}{\uline{linkedin.com/in/uttkarsh-thakur-76979228a}}} \hspace{1pt} $|$
    \hspace{1pt} \mbox{\faGithub \hspace{2pt} \href{https://github.com/uttkarsh-thakur26/}{\uline{github.com/uttkarsh-thakur26}}} \hspace{1pt} $|$
    \hspace{1pt} \mbox{Greater Noida, India}
    \\ \vspace{-3pt}
\end{center}
% <<< end sections/header.tex

%-----------SKILLS-----------
\section{SKILLS}
\begin{itemize}[leftmargin=0in, label={}, itemsep=-1pt]
  \small \item
    \textbf{Languages}: Java, Python, TypeScript, SQL, C++ \\
    \textbf{Backend}: Spring Boot, Spring Web MVC, Spring Data JPA, Hibernate, REST API Design, Bean Validation, FastAPI, Flask, Node.js \\
    \textbf{Databases}: PostgreSQL, MySQL, H2, MongoDB, Flyway Migrations, JPA/Hibernate ORM, Schema Design \\
    \textbf{Frontend}: React.js, TypeScript, Tailwind CSS, Vite, HTML/CSS \\
    \textbf{Testing \& Tools}: JUnit 5, Mockito, AssertJ, Docker, Maven, Git \& GitHub, Swagger/OpenAPI, Postman \\
    \textbf{AI/ML}: RAG Pipelines, Semantic Search, FAISS, LangChain, Transformers, PyTorch, TensorFlow, Scikit-learn
\end{itemize}

%-----------PROJECTS-----------
\section{PROJECTS}
  \resumeSubHeadingListStart

    \resumeProjectHeading
      {\textbf{BillSync} $|$ \textit{Java 21, Spring Boot, Spring Data JPA, PostgreSQL, Flyway, React, TypeScript, Docker}}{2026}
      \resumeItemListStart
        \resumeItem{Built a \textbf{full-stack web application} for splitting group expenses: members log what they paid, and the app works out who owes whom, backed by a \textbf{Java Spring Boot} REST API of \textbf{16 endpoints} and a \textbf{React} and \textbf{TypeScript} interface.}
        \resumeItem{Solved a rounding problem that makes shared balances drift over time, so every split adds up to the \textbf{exact total} with no lost fractions; verified automatically across thousands of amount and group-size combinations.}
        \resumeItem{Developed a \textbf{settle-up feature} that turns a tangle of debts into the fewest payments possible, reducing a 5-member group with 8 outstanding debts to just \textbf{4 transfers}.}
        \resumeItem{Designed the \textbf{PostgreSQL} database of \textbf{6 tables} with version-controlled migrations, and tuned a slow screen that made \textbf{22 database calls} down to \textbf{1}.}
        \resumeItem{Wrote \textbf{73 automated tests} covering business rules and every API endpoint, so new changes ship without breaking existing behaviour.}
      \resumeItemListEnd

    \resumeProjectHeading
      {\textbf{CodeNav AI} $|$ \textit{FastAPI, React.js, LangChain, FAISS, Sentence-Transformers, OpenAI API, GitPython}}{2026}
      \resumeItemListStart
        \resumeItem{Built an \textbf{AI-powered search tool} that lets a developer ask questions about an unfamiliar codebase in plain English and get back the relevant code, using \textbf{large language models} and \textbf{semantic search} (\textbf{RAG}).}
        \resumeItem{Improved answer quality by splitting code along real function and class boundaries instead of fixed-size blocks, so results come back complete and in context.}
        \resumeItem{Cut the time to open a codebase from \textbf{over 20 seconds to under 0.1 seconds}, a \textbf{99.5\% improvement}, by saving the search index to disk instead of rebuilding it each time.}
      \resumeItemListEnd

    \resumeProjectHeading
      {\textbf{Motion Sensing System} $|$ \textit{TensorFlow/Keras, VGG16, Flask, Streamlit, OpenCV, Python}}{2024}
      \resumeItemListStart
        \resumeItem{Built a system that \textbf{recognises human activities from live video}, adapting a pre-trained image-recognition model (\textbf{VGG16}) rather than training from scratch, and delivered it as a live dashboard.}
      \resumeItemListEnd

  \resumeSubHeadingListEnd

%-----------RESEARCH-----------
\section{RESEARCH \& TALKS}
  \resumeSubHeadingListStart
    \resumeProjectHeading
      {\textbf{Agreement Loops in Multi-Agent Systems} $|$ \textit{Seminar Talk, Literature Review}}{2026}
      \resumeItemListStart
        \resumeItem{Reviewed current research on \textbf{how independent AI agents reach agreement} with one another, and presented the findings as a \textbf{seminar talk} covering what works today and what remains unsolved.}
      \resumeItemListEnd
  \resumeSubHeadingListEnd

% >>> inlined from sections/education.tex (generated by tools/flatten.py; edit the original)
%-----------EDUCATION-----------
\section{EDUCATION}
  \resumeSubHeadingListStart
    \resumeSubheading
      {Bennett University}{Aug. 2023 -- May 2027}
      {B.Tech in Computer Science and Engineering $|$ \textbf{CGPA: 8.17/10}}{Greater Noida, India}
      \resumeItemListStart
        \resumeItem{\textbf{Coursework}: Data Structures \& Algorithms, Object-Oriented Programming, Database Management Systems, Operating Systems, Computer Networks, Deep Learning, Neural Networks, System Design}
      \resumeItemListEnd
    \resumeSubheading
      {DAV Public School}{Graduated 2022}
      {Senior Secondary Education (Class XII)}{Ranchi, India}
  \resumeSubHeadingListEnd
% <<< end sections/education.tex

% >>> inlined from sections/certifications.tex (generated by tools/flatten.py; edit the original)
%-----------CERTIFICATIONS-----------
\section{CERTIFICATIONS}
  \resumeSubHeadingListStart
    \resumeSubItem{Build Basic Generative Adversarial Networks (GANs) --- DeepLearning.AI via Coursera \hfill {\color{dark-grey}Apr. 2026}}
    \resumeSubItem{Neural Networks and Deep Learning --- DeepLearning.AI via Coursera \hfill {\color{dark-grey}Apr. 2025}}
    \resumeSubItem{Machine Learning: Clustering \& Retrieval --- University of Washington via Coursera \hfill {\color{dark-grey}Nov. 2024}}
  \resumeSubHeadingListEnd
% <<< end sections/certifications.tex

% >>> inlined from sections/achievements.tex (generated by tools/flatten.py; edit the original)
%-----------ACHIEVEMENTS-----------
\section{ACHIEVEMENTS}
  \resumeSubHeadingListStart
    \resumeSubItem{Ranked \textbf{101\textsuperscript{st} nationally} at \textbf{Smart India Hackathon 2024}}
    \resumeSubItem{\textbf{CodeChef}: \textbf{2-star} rated \quad $|$ \quad \textbf{LeetCode}: Contest rating \textbf{1500+}}
    \resumeSubItem{Competed in \textbf{Hackaccino Hackathon}, Bennett University (2024)}
  \resumeSubHeadingListEnd
% <<< end sections/achievements.tex

\end{document}
